% Zumdahl Chapter 5 -- Gases. ELABORATED notes: fuller exposition, two or
% more worked examples per idea. 5 block lessons.
\documentclass{apchem-notes}

\apcunitword{Chapter}
\apcunit{5}
\apcunittitle{Gases}
\apcdoctitle{Guided Notes}
\apcsubtitle{Zumdahl \S5.1--5.9 \textbullet\ PDF pp.~226--281 \textbullet\ 5 blocks}

\begin{document}
\apctitleblock

\begin{apcnote}[How this chapter fits the AP course]
Most of this chapter you have already met: \S5.2--5.3, 5.5, 5.6, and
5.8 are CED topics \textbf{3.4--3.6}, taught in Unit 3. Read those sections
as \emph{deeper practice}, not new material.

Two things here are genuinely new. \textbf{\S5.4 gas stoichiometry} closes
the loop between Chapter 3's mole ratios and gas volumes --- it is the
reason to work this chapter. And \textbf{collecting a gas over water}
(\S5.5) is a standard laboratory technique that appears on the AP exam and
is covered nowhere else.

\S5.7 (Graham's law of effusion) is marked \textbf{ENRICHMENT} --- it is not
on the current CED. \S5.10 (atmospheric chemistry) is skipped entirely.
\end{apcnote}

% =====================================================================
\blocklesson{Pressure and the Empirical Gas Laws}{Zumdahl \S5.1--5.2}

\begin{objectives}
  \item Explain what gas pressure is and convert among its units.
  \item State Boyle's, Charles's, and Avogadro's laws and use each.
\end{objectives}

\reading{Zumdahl \S5.1--5.2}{227--235}

\begin{warmup}[3]
  \item Convert \SI{25}{\celsius} to kelvin: \answerblank[2.6cm]{\SI{298}{\kelvin}}
  \item Moles in \SI{32.0}{\gram} \ce{O2}: \answerblank[2.6cm]{\SI{1.00}{\mole}}
  \item KMT: average kinetic energy depends only on:
        \answerblank[3cm]{temperature}
\end{warmup}

\segment{INSTRUCTION A}{25}{What pressure is, and how we measure it}

\notesectionz{Pressure}{5.1}
\practice{1}

Pressure is \answerblank[3.9cm]{force per unit area}. For a gas it arises
from countless particle collisions with the container walls --- each
collision is tiny, but there are on the order of $10^{23}$ particles.

Atmospheric pressure is literally the \answerblank[3.5cm]{weight of the
air} above you, which is why it falls with altitude.

\notesub{Torricelli's barometer}

Invented in 1643: a tube of mercury inverted in a dish. Atmospheric pressure
pushing down on the dish holds the column up. At sea level the column
averages \answerblank[2cm]{\SI{760}{\milli\meter}} of mercury.

\begin{apcnote}
The famous crushed-can demonstration is the same physics. Boiling water
fills a can with steam; capping and cooling condenses that steam to a few
drops of liquid, leaving almost nothing inside to push outward. Atmospheric
pressure --- which was always there --- then crushes the can. Nothing
``sucked'' it; the inside simply stopped pushing back.
\end{apcnote}

\notesub{Units you must be able to convert}

\[
  \SI{1}{\atm} = \answerblank[2.2cm]{\SI{760}{\mmHg}}
  = \answerblank[2.0cm]{\SI{760}{\torr}}
  = \SI{101325}{\pascal} = \answerblank[2.6cm]{\SI{101.325}{\kilo\pascal}}
\]

\segment{GUIDED PRACTICE}{15}{Pressure conversions}

\begin{enumerate}[leftmargin=*,itemsep=0.35em]
  \item \SI{2.50}{\atm} in torr: \answerblank[3cm]{\SI{1900}{\torr}}
  \item \SI{380}{\mmHg} in atm: \answerblank[3cm]{\SI{0.500}{\atm}}
  \item \SI{202.65}{\kilo\pascal} in atm: \answerblank[3cm]{\SI{2.000}{\atm}}
\end{enumerate}

\segment{INSTRUCTION B}{20}{Three empirical laws}

\notesectionz{Boyle, Charles, and Avogadro}{5.2}
\practice{5}

Each law fixes two variables and relates the other two:

\renewcommand{\arraystretch}{1.45}
\noindent\begin{tabular}{@{}p{2.6cm}p{2.9cm}p{3.0cm}p{4.6cm}@{}}
\toprule
\textbf{Law} & \textbf{Held constant} & \textbf{Relationship} & \textbf{Working form}\\
\midrule
\term{Boyle} & $n$, $T$ &
  $P$ and $V$ \answerblank[2.2cm]{inversely} &
  \answerblank[2.6cm]{$P_1V_1 = P_2V_2$}\\
\term{Charles} & $n$, $P$ &
  $V$ and $T$ \answerblank[2.2cm]{directly} &
  \answerblank[2.6cm]{$V_1/T_1 = V_2/T_2$}\\
\term{Avogadro} & $P$, $T$ &
  $V$ and $n$ \answerblank[2.2cm]{directly} &
  \answerblank[2.6cm]{$V_1/n_1 = V_2/n_2$}\\
\bottomrule
\end{tabular}

\begin{apctrap}
Charles's and Avogadro's laws are ratios, so temperature \emph{must} be in
kelvin. Celsius has an arbitrary zero: doubling from \SI{10}{\celsius} to
\SI{20}{\celsius} is not doubling the absolute temperature at all --- it is
283 K to 293 K, a 3.5\% rise.
\end{apctrap}

\begin{workedexample}[Worked example 1: Boyle]
A \SI{5.00}{\liter} sample at \SI{1.20}{\atm} is compressed to
\SI{2.00}{\liter} at constant temperature.
\[
  P_2 = \frac{P_1V_1}{V_2} = \frac{(1.20)(5.00)}{2.00} = \SI{3.00}{\atm}
\]
Sanity check: the volume shrank, so the pressure must rise. It did.
\end{workedexample}

\begin{workedexample}[Worked example 2: Charles]
A balloon holds \SI{2.50}{\liter} at \SI{27}{\celsius}. What is its volume
at \SI{127}{\celsius}, pressure constant?
\[
  V_2 = V_1\frac{T_2}{T_1} = 2.50 \times \frac{400.}{300.} = \SI{3.33}{\liter}
\]
Using Celsius would have given $2.50 \times 127/27 = \SI{11.8}{\liter}$ ---
more than three times too large. This is why the conversion matters.
\end{workedexample}

\segment{APPLICATION}{20}{Combining the laws}

Any two states of a fixed gas sample are related by the
\term{combined gas law}:
\[
  \frac{P_1V_1}{T_1} = \answerblank[2.4cm]{$\dfrac{P_2V_2}{T_2}$}
\]

\begin{enumerate}[leftmargin=*,itemsep=0.4em]
  \item A gas occupies \SI{4.00}{\liter} at \SI{2.00}{\atm} and
        \SI{300.}{\kelvin}. Find its volume at \SI{1.00}{\atm} and
        \SI{450.}{\kelvin}. \workspace[2.4cm]
        \keyonly{\begin{apcanswer}$V_2 = 4.00 \times \tfrac{2.00}{1.00}
        \times \tfrac{450.}{300.} = \SI{12.0}{\liter}$\end{apcanswer}}
  \item Explain, using particle collisions, why halving the volume at
        constant temperature doubles the pressure.
        \answerlines[3]{The same number of particles move at the same
        average speed, but the wall area is halved and the distance between
        collisions is shorter --- so each unit of wall area is struck twice
        as often, doubling the force per area.}
\end{enumerate}

\begin{exitticket}
A rigid steel cylinder of gas is heated from \SI{300}{\kelvin} to
\SI{600}{\kelvin}. What happens to its volume and its pressure?
\answerlines[2]{Volume is fixed by the rigid container. Pressure doubles ---
with $V$ and $n$ constant, $P \propto T$ in kelvin.}
\end{exitticket}

% =====================================================================
\blocklesson{The Ideal Gas Law}{Zumdahl \S5.3}

\begin{objectives}
  \item Apply $PV = nRT$ in all rearrangements.
  \item Compute gas density and molar mass from $P$, $T$, and $M$.
\end{objectives}

\reading{Zumdahl \S5.3}{235--241}

\begin{warmup}[4]
  \item \SI{1.50}{\atm} in torr: \answerblank[2.6cm]{\SI{1140}{\torr}}
  \item Boyle's law relates which two variables?
        \answerblank[2.6cm]{$P$ and $V$}
  \item \SI{127}{\celsius} in kelvin: \answerblank[2.4cm]{\SI{400.}{\kelvin}}
  \item Molar mass of \ce{CO2}: \answerblank[3cm]{\SI{44.01}{\gram\per\mole}}
\end{warmup}

\segment{INSTRUCTION A}{25}{One equation containing all three laws}

\notesectionz{$PV = nRT$}{5.3}
\practice{5}

\[
  PV = nRT \qquad
  R = \answerblank[4.2cm]{\SI{0.08206}{\liter\atm\per\mole\per\kelvin}}
\]

$R$ carries units, and those units dictate everything else: pressure in
\answerblank[1.4cm]{atm}, volume in \answerblank[1.4cm]{L},
temperature in \answerblank[1.4cm]{K}. Convert \emph{before} substituting,
never after.

\notesub{Standard temperature and pressure}

STP is \answerblank[2.2cm]{\SI{0}{\celsius}} and
\answerblank[1.6cm]{\SI{1}{\atm}}. At STP one mole of any ideal gas occupies
the \term{molar volume}:
\[
  V = \frac{nRT}{P} = \frac{(1)(0.08206)(273.15)}{1} =
  \answerblank[2.4cm]{\SI{22.42}{\liter}}
\]

\begin{apcnote}
That 22.42 L applies \emph{only} at STP. Students who memorize it and then
use it at \SI{25}{\celsius} get answers about 9\% low. When conditions are
not STP, go back to $PV = nRT$.
\end{apcnote}

\segment{GUIDED PRACTICE}{15}{Rearrangement drill}

\begin{enumerate}[leftmargin=*,itemsep=0.35em]
  \item Solve for $n$: \answerblank[3cm]{$n = PV/RT$}
  \item Volume of \SI{2.50}{\mole} at \SI{1.50}{\atm}, \SI{300.}{\kelvin}:
        \answerblank[3cm]{\SI{41.0}{\liter}}
  \item Pressure of \SI{0.250}{\mole} in \SI{5.00}{\liter} at
        \SI{273}{\kelvin}: \answerblank[3cm]{\SI{1.12}{\atm}}
  \item Moles in \SI{10.0}{\liter} at \SI{2.00}{\atm} and
        \SI{350.}{\kelvin}: \answerblank[3cm]{\SI{0.696}{\mole}}
\end{enumerate}

\segment{INSTRUCTION B}{20}{Density and molar mass}

\notesectionz{Two useful derivations}{5.3}
\practice{5}

Since $n = m/M$, substituting into $PV = nRT$ gives
\[
  d = \frac{m}{V} = \answerblank[2.2cm]{$\dfrac{PM}{RT}$}
  \qquad\text{and}\qquad
  M = \answerblank[2.2cm]{$\dfrac{dRT}{P}$}
\]

You do not need to memorize these if you can derive them --- but you must be
able to get there quickly.

\begin{workedexample}[Worked example 1: density of a known gas]
Find the density of \ce{CO2} at STP.
\[
  d = \frac{PM}{RT} = \frac{(1.00)(44.01)}{(0.08206)(273.15)} = \SI{1.96}{\gram\per\liter}
\]
Equivalently: one mole (\SI{44.01}{\gram}) occupies \SI{22.42}{\liter}, and
$44.01/22.42 = 1.96$. Both routes should agree --- use the second as a
check.
\end{workedexample}

\begin{workedexample}[Worked example 2: identifying an unknown gas]
An unknown gas has a density of \SI{1.34}{\gram\per\liter} at STP.
\[
  M = d \times 22.42 = 1.34 \times 22.42 = \SI{30.0}{\gram\per\mole}
\]
Candidates at 30 g/mol: \ce{NO} (30.01) or \ce{C2H6} (30.07). Density alone
cannot distinguish them --- you would need another property.
\end{workedexample}

\begin{apctrap}
Gas density depends on pressure and temperature, unlike the density of a
solid or liquid. Quoting ``the density of \ce{CO2}'' without stating
conditions is meaningless. Denser gas = higher molar mass \emph{only when
$P$ and $T$ match}.
\end{apctrap}

\segment{APPLICATION}{20}{Working backwards}

\begin{enumerate}[leftmargin=*,itemsep=0.4em]
  \item A \SI{0.500}{\gram} sample of a gas occupies \SI{288}{\milli\liter}
        at \SI{25}{\celsius} and \SI{1.00}{\atm}. Find its molar mass.
        \workspace[2.6cm]
        \keyonly{\begin{apcanswer}$n = PV/RT = (1.00)(0.288)/((0.08206)(298))
        = \SI{0.01178}{\mole}$; $M = 0.500/0.01178 =
        \SI{42.4}{\gram\per\mole}$ (about propane, \ce{C3H8})\end{apcanswer}}
  \item Two identical balloons at the same $T$ and $P$ hold helium and
        argon. Compare their masses and their densities, with reasoning.
        \answerlines[3]{Equal moles (same $P$, $V$, $T$), so argon --- with
        roughly ten times the molar mass --- has both the greater mass and
        the greater density. Density scales with molar mass at fixed
        conditions.}
\end{enumerate}

\begin{exitticket}
Why does a helium balloon rise while an identical balloon of \ce{CO2} sinks?
Answer in terms of density. \answerlines[2]{At the same $T$ and $P$, density
is proportional to molar mass. Helium (4.00) is far less dense than air
($\approx$29), so it is buoyant; \ce{CO2} (44.01) is denser than air and
sinks.}
\end{exitticket}

% =====================================================================
\blocklesson{Gas Stoichiometry}{Zumdahl \S5.4}

\begin{objectives}
  \item Include gas volumes in stoichiometric calculations.
  \item Choose correctly between the molar volume shortcut and $PV = nRT$.
\end{objectives}

\reading{Zumdahl \S5.4}{241--245}

\begin{warmup}[3]
  \item Molar volume at STP: \answerblank[2.6cm]{\SI{22.42}{\liter}}
  \item $d$ of \ce{O2} at STP: \answerblank[3cm]{\SI{1.43}{\gram\per\liter}}
  \item Moles in \SI{50.0}{\gram} \ce{CaCO3} ($M=100.09$):
        \answerblank[2.8cm]{\SI{0.500}{\mole}}
\end{warmup}

\segment{INSTRUCTION A}{25}{Volumes join the mole map}

\notesectionz{The extended road map}{5.4}
\practice{5}

Chapter 3 routed everything through moles. Gases simply add one more
on-ramp and off-ramp:

\begin{center}
\begin{tikzpicture}[node distance=2.3cm, font=\sffamily\small]
  \node[draw=apcaccent, thick, rounded corners, minimum height=0.85cm] (ga) {g A};
  \node[draw=apcaccent, thick, rounded corners, minimum height=0.85cm, right=of ga] (ma) {mol A};
  \node[draw=apcaccent, thick, rounded corners, minimum height=0.85cm, right=of ma] (mb) {mol B};
  \node[draw=apcaccent, thick, rounded corners, minimum height=0.85cm, right=of mb] (vb) {L of gas B};
  \draw[-{Stealth}, thick] (ga) -- node[above,font=\sffamily\scriptsize]{$\div M_A$} (ma);
  \draw[-{Stealth}, thick] (ma) -- node[above,font=\sffamily\scriptsize]{ratio} (mb);
  \draw[-{Stealth}, thick] (mb) -- node[above,font=\sffamily\scriptsize]{$PV{=}nRT$} (vb);
\end{tikzpicture}
\end{center}

Only the middle arrow uses the balanced equation. The last arrow has two
possible tools:

\begin{itemize}[leftmargin=*,itemsep=0.2em]
  \item At STP: multiply by \answerblank[2.4cm]{\SI{22.42}{\liter\per\mole}}
        --- fast.
  \item Any other conditions: use \answerblank[2.2cm]{$V = nRT/P$}.
\end{itemize}

\begin{workedexample}[Worked example 1: Zumdahl's quicklime problem]
Calculate the volume of \ce{CO2} at STP produced by decomposing
\SI{152}{\gram} of \ce{CaCO3}:
\[ \ce{CaCO3(s) -> CaO(s) + CO2(g)} \]
\begin{align*}
  n(\ce{CaCO3}) &= \frac{152}{100.09} = \SI{1.52}{\mole}\\
  n(\ce{CO2}) &= \SI{1.52}{\mole} \quad (1{:}1)\\
  V &= 1.52 \times 22.42 = \SI{34.1}{\liter}
\end{align*}
\end{workedexample}

\begin{workedexample}[Worked example 2: not at STP]
What volume of \ce{O2} at \SI{25}{\celsius} and \SI{1.00}{\atm} is needed to
burn \SI{10.0}{\gram} of \ce{CH4}?
\[ \ce{CH4 + 2 O2 -> CO2 + 2 H2O} \]
\begin{align*}
  n(\ce{CH4}) &= \frac{10.0}{16.04} = \SI{0.623}{\mole}\\
  n(\ce{O2}) &= 2(0.623) = \SI{1.25}{\mole}\\
  V &= \frac{nRT}{P} = \frac{(1.25)(0.08206)(298)}{1.00} = \SI{30.5}{\liter}
\end{align*}
The molar-volume shortcut would have given \SI{28.0}{\liter} --- wrong by
8\%, because these are not standard conditions.
\end{workedexample}

\segment{GUIDED PRACTICE}{15}{Gas stoichiometry drill}

\begin{enumerate}[leftmargin=*,itemsep=0.4em]
  \item \ce{2 KClO3(s) -> 2 KCl(s) + 3 O2(g)}. What volume of \ce{O2} at
        STP comes from \SI{12.25}{\gram} \ce{KClO3} ($M = 122.55$)?
        \workspace[2.2cm]
        \keyonly{\begin{apcanswer}$n = 0.1000$ mol $\times \tfrac32 =
        0.1500$ mol \ce{O2} $\times 22.42 =
        \SI{3.36}{\liter}$\end{apcanswer}}
  \item \ce{N2(g) + 3 H2(g) -> 2 NH3(g)}. All gases at the same $T$ and
        $P$: what volume of \ce{H2} reacts with \SI{5.0}{\liter} of
        \ce{N2}? \answerblank[3cm]{\SI{15}{\liter}}
\end{enumerate}

\begin{apcnote}
That second problem needed no molar masses at all. At fixed $T$ and $P$,
volume is proportional to moles (Avogadro), so gas volumes combine in the
same ratio as the coefficients. This shortcut works \emph{only} when every
substance involved is a gas at the same conditions.
\end{apcnote}

\segment{APPLICATION}{20}{Limiting reactant with gases}

\SI{4.00}{\liter} of \ce{H2} and \SI{4.00}{\liter} of \ce{O2}, both at
\SI{1.00}{\atm} and \SI{273}{\kelvin}, react to form water.

\begin{enumerate}[leftmargin=*,itemsep=0.4em]
  \item Write the balanced equation.
        \answerblank[5cm]{\ce{2 H2 + O2 -> 2 H2O}}
  \item Identify the limiting reactant using volumes directly.
        \answerlines[2]{The ratio needed is 2 \ce{H2} : 1 \ce{O2}. Equal
        volumes means equal moles, so hydrogen runs out first --- \ce{H2} is
        limiting.}
  \item What volume of \ce{O2} remains unreacted?
        \answerblank[3cm]{\SI{2.00}{\liter}}
\end{enumerate}

\begin{exitticket}
Why can you compare gas volumes directly to find a limiting reactant, but
never compare gas masses that way? \answerlines[2]{At fixed $T$ and $P$,
volume is proportional to moles, so volume ratios equal mole ratios. Mass
ratios do not, because the gases have different molar masses.}
\end{exitticket}

% =====================================================================
\blocklesson{Partial Pressures and Gas Collection}{Zumdahl \S5.5}

\begin{objectives}
  \item Apply Dalton's law and mole fractions.
  \item Correct for water vapor when a gas is collected over water.
\end{objectives}

\reading{Zumdahl \S5.5}{245--251}

\begin{warmup}[3]
  \item Volume of \SI{0.500}{\mole} at STP:
        \answerblank[2.8cm]{\SI{11.2}{\liter}}
  \item \ce{N2 + 3 H2 -> 2 NH3}: \SI{6}{\liter} \ce{H2} needs how much
        \ce{N2}? \answerblank[2.4cm]{\SI{2}{\liter}}
  \item Molar mass from $d$ at STP: \answerblank[3.4cm]{$M = d \times 22.42$}
\end{warmup}

\segment{INSTRUCTION A}{25}{Dalton's law}

\notesectionz{Each gas acts independently}{5.5}
\practice{5}

In an ideal gas, particles do not interact --- so each component exerts
pressure as if the others were absent:
\[
  P_{\text{total}} = \answerblank[4cm]{$P_1 + P_2 + P_3 + \cdots$}
\]

Zumdahl derives the mole-fraction relationship directly from $n = P(V/RT)$:
since each $n_i$ is proportional to $P_i$ with the same constant,
\[
  \chi_1 = \frac{n_1}{n_{\text{total}}} = \answerblank[2.6cm]{$\dfrac{P_1}{P_{\text{total}}}$}
  \qquad\Rightarrow\qquad
  P_1 = \answerblank[2.6cm]{$\chi_1 P_{\text{total}}$}
\]

\begin{apcnote}
The identity of the gases is irrelevant. Three moles of helium and three
moles of \ce{SF6} --- molar masses 4 and 146 --- contribute exactly the same
partial pressure at the same $T$ and $V$. Pressure counts
\emph{particles}, not mass.
\end{apcnote}

\begin{workedexample}[Worked example 1]
A vessel holds \SI{0.30}{\mole} \ce{N2}, \SI{0.20}{\mole} \ce{O2}, and
\SI{0.50}{\mole} \ce{He} at a total pressure of \SI{2.00}{\atm}.
\[
  \chi_{\ce{N2}} = \frac{0.30}{1.00} = 0.30
  \qquad P_{\ce{N2}} = 0.30 \times 2.00 = \SI{0.60}{\atm}
\]
Check: the three partial pressures (0.60, 0.40, 1.00) sum to 2.00. Always
verify this.
\end{workedexample}

\segment{GUIDED PRACTICE}{15}{Partial pressure drill}

A \SI{10.0}{\liter} flask at \SI{300.}{\kelvin} contains \SI{0.400}{\mole}
\ce{Ar} and \SI{0.600}{\mole} \ce{Ne}.
\begin{enumerate}[leftmargin=*,itemsep=0.35em]
  \item Total pressure: \answerblank[3cm]{\SI{2.46}{\atm}}
  \item $P_{\ce{Ar}}$: \answerblank[3cm]{\SI{0.985}{\atm}}
  \item If all the Ne escaped, $P_{\ce{Ar}}$ would be:
        \answerblank[4.4cm]{unchanged, \SI{0.985}{\atm}}
\end{enumerate}

\segment{INSTRUCTION B}{20}{Collecting a gas over water}

\notesectionz{The laboratory correction}{5.5}
\practice{4}

A gas produced in a reaction is often collected by displacing water from an
inverted container. The collected sample is then a
\answerblank[2.2cm]{mixture}: the gas you made, plus
\answerblank[2.6cm]{water vapour}.

\[
  P_{\text{total}} = P_{\text{gas}} + \answerblank[2.4cm]{$P_{\ce{H2O}}$}
  \qquad\Rightarrow\qquad
  P_{\text{gas}} = P_{\text{total}} - P_{\ce{H2O}}
\]

The vapour pressure of water depends only on
\answerblank[2.7cm]{temperature} and is read from a table:

\renewcommand{\arraystretch}{1.25}
\begin{center}
\begin{tabular}{@{}lccccc@{}}
\toprule
$T$ (\si{\celsius}) & 15 & 20 & 25 & 30 & 35\\
$P_{\ce{H2O}}$ (torr) & 12.8 & 17.5 & 23.8 & 31.8 & 42.2\\
\bottomrule
\end{tabular}
\end{center}

\begin{workedexample}[Worked example 2: hydrogen over water]
Hydrogen is collected over water at \SI{25}{\celsius}. The total pressure is
\SI{755}{\torr} and the volume is \SI{0.250}{\liter}. How many moles of
\ce{H2} were collected?
\begin{align*}
  P_{\ce{H2}} &= 755 - 23.8 = \SI{731}{\torr} = \SI{0.962}{\atm}\\
  n &= \frac{PV}{RT} = \frac{(0.962)(0.250)}{(0.08206)(298)} = \SI{9.83e-3}{\mole}
\end{align*}
Forgetting the correction would have overstated the answer by about 3\%.
\end{workedexample}

\begin{apctrap}
The correction is always a \emph{subtraction} --- the water vapour is
inflating your reading. Students who add it, or who skip it entirely, report
too much gas. Any AP question mentioning ``collected over water'' is
signalling that this step is required.
\end{apctrap}

\segment{APPLICATION}{20}{Full collection analysis}

Oxygen from the decomposition of \ce{KClO3} is collected over water at
\SI{20}{\celsius}. The total pressure is \SI{742}{\torr}; the volume is
\SI{0.500}{\liter}.

\begin{enumerate}[leftmargin=*,itemsep=0.4em]
  \item Partial pressure of \ce{O2} in torr and atm:
        \answerblank[5.5cm]{\SI{724.5}{\torr} $=$ \SI{0.9539}{\atm}}
  \item Moles of \ce{O2} collected: \workspace[2.2cm]
        \keyonly{\begin{apcanswer}$n = (0.9539)(0.500)/((0.08206)(293)) =
        \SI{0.01983}{\mole}$\end{apcanswer}}
  \item Mass of \ce{KClO3} that decomposed, given
        \ce{2 KClO3 -> 2 KCl + 3 O2} ($M = 122.55$). \workspace[2.2cm]
        \keyonly{\begin{apcanswer}$n(\ce{KClO3}) = \tfrac23(0.01983) =
        \SI{0.01322}{\mole}$; $m = 0.01322 \times 122.55 =
        \SI{1.62}{\gram}$\end{apcanswer}}
\end{enumerate}

\begin{exitticket}
Why does the vapour pressure of water in a collection tube depend on
temperature but not on which gas is being collected?
\answerlines[2]{It is set by the equilibrium between liquid water and its
own vapour, which depends only on temperature. The other gas exerts its
pressure independently (Dalton's law).}
\end{exitticket}

% =====================================================================
\blocklesson{Kinetic Molecular Theory and Real Gases}{Zumdahl \S5.6--5.9}

\begin{objectives}
  \item State the KMT postulates and connect them to the gas laws.
  \item Predict and explain deviations from ideal behaviour.
\end{objectives}

\reading{Zumdahl \S5.6--5.9}{251--264}

\begin{warmup}[4]
  \item $P_{\ce{H2O}}$ at \SI{25}{\celsius}: \answerblank[2.4cm]{\SI{23.8}{\torr}}
  \item $P_A$ if $\chi_A = 0.25$, $P_{\text{tot}} = \SI{4.0}{\atm}$:
        \answerblank[2.4cm]{\SI{1.0}{\atm}}
  \item Volume of \SI{2.00}{\mole} at STP:
        \answerblank[2.8cm]{\SI{44.8}{\liter}}
  \item Gas collected over water: add or subtract $P_{\ce{H2O}}$?
        \answerblank[2.4cm]{subtract}
\end{warmup}

\segment{INSTRUCTION A}{25}{The model}

\notesectionz{KMT postulates}{5.6}
\practice{1}

An ideal gas consists of particles that:
\begin{enumerate}[leftmargin=*,itemsep=0.15em]
  \item are in constant \answerblank[2.2cm]{random} motion;
  \item have \answerblank[3.7cm]{negligible volume} relative to the container;
  \item exert \answerblank[4.0cm]{no attractive forces} on one another;
  \item collide \answerblank[2.4cm]{elastically};
  \item have average kinetic energy proportional to
        \answerblank[4.5cm]{absolute temperature}.
\end{enumerate}

\[
  KE_{\text{avg}} = \tfrac32 RT \ \text{per mole}
  \qquad
  u_{\text{rms}} = \answerblank[2.6cm]{$\sqrt{\dfrac{3RT}{M}}$}
\]

\begin{apcnote}
Read the speed equation carefully: at a given temperature, every gas has the
\emph{same} average kinetic energy, but speed depends on molar mass. Heavier
molecules must move more slowly to carry the same energy --- speed scales as
$1/\sqrt{M}$.
\end{apcnote}

\segment{GUIDED PRACTICE}{15}{Explaining laws with the model}

\begin{enumerate}[leftmargin=*,itemsep=0.35em]
  \item Why does pressure rise when temperature rises at constant volume?
        \answerlines[2]{Higher $T$ raises average speed, so collisions with
        the walls are both more frequent and more forceful.}
  \item Rank average speed at the same $T$: \ce{He}, \ce{N2}, \ce{CO2}.
        \answerblank[5cm]{\ce{He} $>$ \ce{N2} $>$ \ce{CO2}}
\end{enumerate}

\begin{apcnote}[ENRICHMENT --- beyond the CED]
\textbf{Graham's law of effusion} (\S5.7). Effusion is escape through a tiny
hole. Graham found the rate varies inversely with the square root of molar
mass:
\[
  \frac{\text{rate}_1}{\text{rate}_2} = \sqrt{\frac{M_2}{M_1}}
\]
Zumdahl's example: \ce{H2} (2.016) versus \ce{UF6} (352.02) gives
$\sqrt{352.02/2.016} = 13.2$ --- hydrogen effuses about 13 times faster.
This is not a current CED topic, though the underlying idea
(lighter $\Rightarrow$ faster) certainly is.
\end{apcnote}

\segment{INSTRUCTION B}{20}{Where the ideal model fails}

\notesectionz{Real gases}{5.8--5.9}
\practice{6}

Two postulates break under stress, and they push $PV/nRT$ in
\emph{opposite} directions:

\renewcommand{\arraystretch}{1.4}
\noindent\begin{tabular}{@{}p{3.6cm}p{3.6cm}p{6.0cm}@{}}
\toprule
\textbf{Failed postulate} & \textbf{Conditions} & \textbf{Effect}\\
\midrule
Particles have no volume & \answerblank[2.6cm]{high pressure} &
  gas cannot compress as predicted; $PV/nRT$
  \answerblank[1.4cm]{$>1$}\\
No attractive forces & \answerblank[3.5cm]{low temperature} &
  attractions soften wall collisions, lowering $P$; $PV/nRT$
  \answerblank[1.4cm]{$<1$}\\
\bottomrule
\end{tabular}

Van der Waals corrects both:
\[
  \left(P + \frac{an^2}{V^2}\right)(V - nb) = nRT
\]
where $a$ corrects for \answerblank[2.4cm]{attractions} and $b$ for
\answerblank[3.3cm]{particle volume}. You will not compute with this, but
you should be able to say which term fixes which failure.

\segment{APPLICATION}{20}{Deviation reasoning}

\begin{enumerate}[leftmargin=*,itemsep=0.4em]
  \item Rank by expected deviation from ideality at the same conditions:
        \ce{He}, \ce{N2}, \ce{H2O}. Justify.
        \answerlines[3]{\ce{H2O} $>$ \ce{N2} $>$ \ce{He}. Water hydrogen
        bonds strongly, so its attractive correction is large; \ce{N2} has
        moderate dispersion forces; helium is tiny with almost no
        attractions.}
  \item A gas shows $PV/nRT = 0.95$ at moderate pressure and
        \SI{200}{\kelvin}. Which postulate is failing, and what happens if
        the sample is warmed to \SI{500}{\kelvin}?
        \answerlines[3]{Attractions dominate, lowering the measured
        pressure below ideal. Warming increases kinetic energy so molecules
        overcome those attractions --- the ratio rises back toward 1.}
\end{enumerate}

\begin{exitticket}
State the two conditions under which any real gas behaves most ideally, and
give the reason for each. \answerlines[2]{High temperature --- fast
molecules overcome intermolecular attractions. Low pressure --- particles
are far apart, so their own volume is negligible compared with the
container.}
\end{exitticket}

\end{document}
